% !TeX root = ../main.tex

\xchapter{相关理论与技术}{Related Theories and Techniques}

本章介绍本文研究所涉及的相关理论与关键技术，为后续方法设计提供必要的理论基础。首先介绍惯性测量单元（IMU）与WiFi信道状态信息（CSI）两类传感器的基本感知原理。随后介绍Transformer与Mamba两种深度学习时序建模网络结构，并分析其在长时序数据建模中的特点与适用场景。在此基础上，进一步介绍多模态特征对齐与融合方法，包括对比学习与交叉注意力机制，为后续多模态人体活动识别模型的设计提供理论依据。最后介绍自回归生成机制的基本思想，为人体姿态重构与预测方法的实现奠定基础。

\xsection{传感器}{Sensor}

\xsubsection{惯性测量单元}{Inertial Measurement Unit}

惯性测量单元（Inertial Measurement Unit, IMU）是人体运动感知系统中常用的可穿戴传感器，通常由三轴加速度计和三轴陀螺仪组成，用于测量载体在三维空间中的线性加速度和角速度信息。由于IMU可以固定在人体肢体上并随肢体运动，因此能够连续记录人体运动过程中的动力学变化，在人体活动识别和姿态建模等任务中得到广泛应用。

从观测方式来看，IMU传感器随人体肢体运动实时采集信号，其观测方式可以视为一种基于载体的本体感知方式。通过对加速度和角速度信号进行连续采样，可以获得人体运动过程中各肢体的时序运动信息。

在加速度测量方面，加速度计输出的是载体坐标系下的比力（Specific Force），即线性加速度与重力加速度共同作用的结果。设时刻 $t$ 的加速度计输出为 $a_t$，其观测模型可以表示为：

\begin{equation}
a_t = C_n^b (a_n - g_n) + b_a + n_a
\end{equation}

其中，$C_n^b$ 表示从导航坐标系（Navigation Frame）到载体坐标系（Body Frame）的旋转矩阵，$a_n$ 为人体肢体在世界坐标系下的线性加速度，$g_n$ 为重力加速度向量，$b_a$ 为加速度计零偏误差，$n_a$ 为测量噪声。该模型表明IMU采集的加速度信号同时包含人体运动产生的加速度分量以及重力分量。

陀螺仪用于测量载体的角速度信息，其输出模型可以表示为：

\begin{equation}
\omega_t = \omega_{true} + b_g + n_g
\end{equation}

其中，$\omega_{true}$ 为真实角速度，$b_g$ 为陀螺仪零偏误差，$n_g$ 为测量噪声。角速度信号能够反映人体肢体在运动过程中的旋转变化。

在人体感知任务中，通过在人体关键部位佩戴多个IMU传感器，可以获得人体各肢体的时序运动信号。与视觉传感器相比，IMU能够以较高采样率连续记录人体运动过程中的动力学变化，对人体动作中的微小运动模式具有较高的敏感性，因此在基于可穿戴设备的人体活动识别与人体姿态建模研究中具有重要应用价值。

\xsubsection{WiFi信道状态信息}{WiFi Channel State Information}

WiFi感知技术利用无线信号在传播过程中受到人体运动影响而产生的信道变化，从而实现对人体行为的感知。与IMU等可穿戴传感器不同，WiFi感知属于一种基于环境观测的非接触式感知方式，即通过监测固定空间区域内无线信号传播特性的变化来推断人体运动状态。

现代WiFi系统通常采用正交频分复用（Orthogonal Frequency Division Multiplexing, OFDM）技术进行数据传输。在OFDM系统中，高速数据流被划分为多个子载波并行传输，从而能够有效应对复杂无线信道带来的频率选择性衰落。信道状态信息（Channel State Information, CSI）是物理层对无线信号在发射端与接收端之间传播特性的细粒度刻画，用于描述不同子载波上的信道频率响应。

在频域上，对于第 $k$ 个子载波，接收信号 $Y_k$ 与发送信号 $X_k$ 之间的关系可以表示为：

\begin{equation}
Y_k = H_k X_k + N_k
\end{equation}

其中，$H_k$ 为该子载波上的复数信道增益，即CSI；$N_k$ 为高斯白噪声。CSI通常以复数形式表示，可以分解为幅度和相位：

\begin{equation}
H_k = |H_k| e^{j\angle H_k}
\end{equation}

其中，幅度 $|H_k|$ 表示信号在该子载波上的能量衰减，而相位 $\angle H_k$ 则反映了信号传播路径的时间延迟信息。与仅提供单一信号强度信息的接收信号强度指示（RSSI）相比，CSI能够提供更加丰富的无线传播特征，因此在无线感知研究中得到了广泛应用。

在室内环境中，无线信号通常通过多条路径传播，包括来自墙壁、地面、家具以及人体的反射路径。根据多径传播模型，时变信道频率响应可以表示为：

\begin{equation}
H(f,t) = \sum_{p=1}^{P} \alpha_p(t) e^{-j2\pi f \tau_p(t)}
\end{equation}

其中，$P$ 表示传播路径数量，$\alpha_p(t)$ 为第 $p$ 条路径的衰减系数，$\tau_p(t)$ 为对应的传播时延。当人体在WiFi覆盖区域内运动时，会改变部分传播路径的长度和反射特性，从而引起 $\alpha_p(t)$ 与 $\tau_p(t)$ 的变化。这种变化会进一步反映在CSI幅度和相位的时序变化中，使得CSI信号能够反映人体运动引起的无线传播环境变化。

因此，通过对CSI信号进行时序建模，可以从无线信号的变化中提取人体运动相关特征，从而实现人体行为识别和环境感知等任务。

\xsection{深度学习网络}{Tables}

\xsubsection{Transformer}{Transformer}

Transformer是一种基于注意力机制的序列建模网络结构，在自然语言处理和时序建模任务中得到广泛应用。与循环神经网络（RNN）依赖递归更新隐藏状态不同，Transformer通过自注意力机制（Self-Attention Mechanism）直接建模序列中任意位置之间的依赖关系，从而能够有效捕捉长时序数据中的全局关联信息。在人体运动建模任务中，这种全局建模能力有助于学习人体动作在时间维度上的动态变化模式。

Transformer的核心计算单元为缩放点积注意力（Scaled Dot-Product Attention）。给定输入序列的特征表示，模型首先通过线性映射将其转换为查询向量（Query, $Q$）、键向量（Key, $K$）和值向量（Value, $V$）。注意力输出可以表示为：

\begin{equation}
\text{Attention}(Q, K, V) =
\text{softmax}\left(\frac{QK^T}{\sqrt{d_k}}\right)V
\end{equation}

其中 $d_k$ 为键向量的维度，缩放因子 $\sqrt{d_k}$ 用于稳定训练过程。该计算过程通过度量查询向量与键向量之间的相似度，获得不同时间步之间的注意力权重，从而实现对序列特征的加权聚合。

为了增强模型对不同特征子空间的表达能力，Transformer进一步引入多头注意力机制（Multi-Head Attention）。该机制将输入特征投影到多个子空间中并行计算注意力，再将各个注意力头的结果进行拼接与线性变换，其计算形式如下：

\begin{equation}
\text{MultiHead}(Q, K, V) =
\text{Concat}(\text{head}_1, ..., \text{head}_h)W^O
\end{equation}

\begin{equation}
\text{head}_i =
\text{Attention}(QW_i^Q, KW_i^K, VW_i^V)
\end{equation}

其中 $W_i^Q$、$W_i^K$、$W_i^V$ 和 $W^O$ 为可学习参数矩阵。通过多头注意力结构，模型能够在不同表示子空间中学习序列特征之间的关联关系。

由于自注意力机制本身不包含序列顺序信息，因此Transformer通过位置编码（Positional Encoding）引入时序位置信息。常用的位置编码方式采用正弦与余弦函数生成：

\begin{equation}
PE_{(pos, 2i)} =
\sin(pos / 10000^{2i/d_{model}})
\end{equation}

\begin{equation}
PE_{(pos, 2i+1)} =
\cos(pos / 10000^{2i/d_{model}})
\end{equation}

通过将位置编码与输入特征相加，模型能够感知序列中的时间顺序关系，从而实现对时序数据的有效建模。

Transformer模型结构如图2-1(a)所示。该结构通过多层注意力计算与前馈网络的组合，实现了对序列数据中长距离依赖关系的建模能力。在人体感知任务中，Transformer能够对长时间范围内的人体运动信号进行建模，因此在人体姿态重构等需要高精度时序建模的任务中具有良好的应用效果。

\begin{figure}[h]
    \begin{subfigure}[b]{0.49\linewidth}
        \centering
        \includegraphics[width=\textwidth]{transformer.png}
        \subcaption{Transformer}
    \end{subfigure}
    \begin{subfigure}[b]{0.49\linewidth}
        \centering
        \includegraphics[width=\textwidth]{mamba.png}
        \subcaption{Mamba}
        \label{subfigure:icon}
    \end{subfigure}
    \caption{模型架构图}
    \label{figure:icon}
\end{figure}

\xsubsection{Mamba}{Mamba}

尽管Transformer在长序列建模任务中表现出良好的性能，但其自注意力机制的计算复杂度随序列长度呈二次增长。当处理高采样率传感器产生的长时序数据时，Transformer在计算开销和显存占用方面可能面临一定挑战。为提高长序列建模效率，近年来研究者提出了基于状态空间模型（State Space Model, SSM）的序列建模方法，其中Mamba是一种具有代表性的结构。

状态空间模型通过引入隐状态来描述输入序列的动态演化过程。在离散时间形式下，其基本表达形式可以表示为：

\begin{equation}
h_t = A h_{t-1} + B x_t
\end{equation}

\begin{equation}
y_t = C h_t
\end{equation}

其中，$x_t$ 表示输入序列，$h_t$ 表示系统隐状态，$y_t$ 为输出序列，$A$、$B$ 和 $C$ 为模型参数矩阵。通过递归更新隐状态，状态空间模型能够对长时间范围内的序列信息进行建模。

在传统SSM中，模型参数通常是固定的线性变换，这限制了模型对复杂序列模式的表达能力。Mamba通过引入选择性机制（Selective Mechanism），使部分模型参数能够根据输入动态调整，从而提升模型对不同时间步输入信号的适应能力。具体而言，Mamba将输入序列映射到一组输入相关的参数，用于控制信息在状态更新过程中的保留与遗忘，从而实现对重要特征的动态选择。

与基于自注意力机制的模型相比，Mamba在序列建模过程中采用线性复杂度的计算方式，因此在处理长序列数据时具有更高的计算效率。同时，通过状态空间结构对序列信息进行递归建模，模型能够在保持较低计算开销的情况下捕捉长时依赖关系。

Mamba模型结构如图2-1(b)所示。由于人体活动识别任务通常涉及较长时间范围的传感器数据序列，Mamba在保证建模能力的同时能够降低计算复杂度，因此适用于对长时序人体运动数据进行特征建模。

\xsection{多模态特征融合与对齐}{Figures}

在人体感知任务中，不同类型的传感器通常能够从不同角度反映人体运动信息。例如，惯性测量单元（IMU）能够直接记录人体肢体运动过程中的加速度和角速度变化，反映人体运动的动力学特征；而基于WiFi信号的感知方法则通过分析信道状态信息（CSI）的变化，捕捉人体运动对无线传播环境产生的影响，从而获得环境层面的感知信息。相比单一模态传感器，多模态传感数据能够从不同维度描述人体行为状态，为人体活动识别任务提供更加丰富的信息。

然而，不同模态传感数据在物理属性、数据分布以及特征表达方式上通常存在明显差异，直接进行特征拼接或简单融合往往难以充分利用多模态信息。一方面，模态之间的分布差异可能导致特征空间不一致，从而影响模型的表征能力；另一方面，不同模态所包含的信息具有一定互补性，如果能够有效建立模态之间的关联关系，则有助于提升模型对人体行为的表达能力。

因此，在多模态人体活动识别任务中，需要同时解决模态对齐与模态融合两个问题。模态对齐旨在将来自不同传感器的特征映射到统一的语义空间，从而减少模态之间的分布差异，使模型能够学习到更加通用的表征；而模态融合则通过建立不同模态特征之间的交互关系，充分利用各模态所包含的互补信息，从而提升模型的整体感知能力。

基于上述考虑，本节将介绍实现多模态特征对齐与融合的相关方法，包括用于学习跨模态共享表示的对比学习方法，以及用于建模不同模态特征交互关系的交叉注意力机制。

\xsubsection{对比学习}{Constrictive Learning}

在多模态人体活动识别任务中，惯性测量单元（IMU）记录的动力学特征与WiFi记录的空间场特征在物理属性、数据分布及维度上存在本质差异，形成了严峻的模态鸿沟。传统的融合方法往往假设异构特征在原始空间或浅层映射后即可线性叠加，忽略了特征分布的不一致性，容易导致多模态信息的相互干扰而非互补。为了解决这一问题，本研究引入跨模态对比学习（Cross-Modal Contrastive Learning）作为表征对齐的关键机制。其核心思想并非直接对异构数据进行分类，而是通过在潜在语义空间中度量样本间的相似性，强制语义相关的配对样本（正样本）相互靠近，而语义无关的非配对样本（负样本）相互远离，从而构建一个共享的、具备语义一致性的高维流形。

\begin{figure}[h]
    \centering
    \includegraphics[width=.6\textwidth]{对比学习原理示意图.png}
    \caption{对比学习原理示意图}
\end{figure}

受大规模图文预训练模型CLIP的启发，本研究摒弃了传统对比学习中繁琐的二元组或三元组正负样本挖掘过程，转而采用基于批次（Batch-wise）的全局相似度计算范式。假设在一个训练批次中包含 $N$ 个时间同步的传感器数据片段对 $\{(x_i^{I}, x_i^{W})\}_{i=1}^N$，其中 $x_i^{I}$ 和 $x_i^{W}$ 分别代表第 $i$ 个样本的IMU数据与WiFi数据。首先，通过各自的特征提取网络将原始数据映射为高维特征向量，并经过非线性投影头映射至同一维度的度量空间，得到归一化的嵌入向量 $u_i$ 和 $v_i$。在此基础上，算法构建一个 $N \times N$ 的全量相似度矩阵 $S$，其中第 $i$ 行第 $j$ 列的元素 $S_{ij}$ 表示IMU样本 $i$ 与WiFi样本 $j$ 之间的余弦相似度：
\begin{equation}
S_{ij} = \frac{f_{IMU}(x_i^{I}) \cdot f_{WiFi}(x_j^{W})}{\|f_{IMU}(x_i^{I})\| \|f_{WiFi}(x_j^{W})\|} = u_i^T v_j
\label{eq:similarity_matrix}
\end{equation}
在该相似度矩阵中，对角线元素 $S_{ii}$ 对应于时间同步的真实配对样本（即同一时刻的人体动作），应被最大化；而非对角线元素 $S_{ij} (i \neq j)$ 对应于批次内其他时刻或不同动作的样本组合，属于负样本，应被最小化。这一优化目标等价于在批次内进行 $N$ 类分类问题。为了同时利用IMU到WiFi和WiFi到IMU的双向信息，本研究采用对称的对比损失函数。具体而言，对于每一个IMU样本 $i$，将其视为“查询”，计算其与批次内所有WiFi样本的匹配概率，构建IMU-to-WiFi的损失 $L_i^{I \rightarrow W}$：
\begin{equation}
L_i^{I \rightarrow W} = -\log \frac{\exp(S_{ii} / \tau)}{\sum_{k=1}^N \exp(S_{ik} / \tau)}
\label{eq:loss_i2w}
\end{equation}
同理，计算WiFi-to-IMU的损失 $L_i^{W \rightarrow I}$：
\begin{equation}
L_i^{W \rightarrow I} = -\log \frac{\exp(S_{ii} / \tau)}{\sum_{k=1}^N \exp(S_{ki} / \tau)}
\label{eq:loss_w2i}
\end{equation}
其中，$\tau$ 为可学习的温度系数，用于调节Softmax分布的平滑程度，控制模型对困难负样本的挖掘力度。最终的总损失函数为两者的平均值：
\begin{equation}
L_{CL} = \frac{1}{2N} \sum_{i=1}^N (L_i^{I \rightarrow W} + L_i^{W \rightarrow I})
\label{eq:total_cl_loss}
\end{equation}
通过最小化该损失函数，模型在反向传播过程中能够实现两个关键特性：一是“对齐性”（Alignment），即拉近同一动作在异构模态下的特征距离，消除物理属性差异带来的隔阂；二是“均匀性”（Uniformity），即迫使特征向量在超球面上尽可能均匀分布，保留更多的语义信息并防止模式崩塌。这种基于CLIP范式的对比学习策略，为后续的多任务分类网络提供了鲁棒且对齐的特征底座，是解决异构感知数据深度融合的数学基础。

\xsubsection{交叉注意力机制}{Cross-Attention Mechanism}

尽管跨模态对比学习能够将异构数据映射至统一的语义流形，实现了全局特征的对齐，但它并未解决模态间细粒度特征的交互与融合问题。在多模态人体活动识别中，惯性数据与射频数据虽然描述同一动作，但其关注点截然不同：IMU聚焦于肢体的微观动力学变化，而WiFi聚焦于宏观的空间环境扰动。为了有效地融合这两种包含不同语义层级且在时序上可能存在非线性延迟的信息，本研究引入交叉注意力机制（Cross-Attention Mechanism）。与自注意力机制不同，交叉注意力的核心在于其查询向量（Query）与键值对向量（Key-Value）来源于不同的模态，这使得模型能够动态地利用一个模态的信息去通过注意力权重“检索”并聚合另一个模态中的互补特征。

% \begin{figure}[h]
%     \centering
%     \includegraphics[width=\textwidth]{交叉注意力机制原理示意图.png}
%     \caption{交叉注意力机制原理示意图}
% \end{figure}

假设经过特征提取后的IMU特征序列为 $X_A \in \mathbb{R}^{L_A \times d}$，WiFi特征序列为 $X_B \in \mathbb{R}^{L_B \times d}$。以IMU特征为主模态去融合WiFi特征为例，首先通过线性投影生成查询、键和值矩阵。此时，查询矩阵 $Q_A$ 源自IMU特征，而键矩阵 $K_B$ 和值矩阵 $V_B$ 则源自WiFi特征：
\begin{equation}
Q_A = X_A W^Q, \quad K_B = X_B W^K, \quad V_B = X_B W^V
\label{eq:cross_projection}
\end{equation}
其中，$W^Q, W^K, W^V \in \mathbb{R}^{d \times d_k}$ 为可学习的线性投影权重。交叉注意力的计算过程旨在衡量IMU序列中第 $i$ 个时间步的特征与WiFi序列中所有时间步特征的相关性。注意力权重矩阵及最终的融合输出 $O_{A \leftarrow B}$ 计算如下：
\begin{equation}
O_{A \leftarrow B} = \text{softmax}\left(\frac{Q_A K_B^T}{\sqrt{d_k}}\right) V_B
\label{eq:cross_attention_calc}
\end{equation}
在该公式中，点积项 $Q_A K_B^T$ 构建了一个跨模态的相似度图（Affinity Map）。对于IMU中的每一个动作片段，该机制会自动寻找WiFi信号中与之最相关的空间特征片段，并赋予其更高的权重。随后，通过对值矩阵 $V_B$ 的加权求和，将相关的WiFi环境信息注入到IMU的动力学特征中。

这种交互机制在物理层面上具有明确的解释性：它允许模型在处理IMU的加速度突变时，参考同期的WiFi多径变化情况，从而判断该突变是由人体运动引起还是环境噪声导致。为了实现双向的信息互补，通常会并行构建 $O_{A \leftarrow B}$ 与 $O_{B \leftarrow A}$ 两个方向的交叉注意力流，通过多头（Multi-Head）并行计算捕获不同子空间下的模态关联。最终，交叉注意力机制有效地解决了异构数据在时空维度上的局部错位问题，为后续的多任务分类提供了富含多模态上下文的增强表征。

\xsection{自回归生成机制}{Autoregressive Mechanism}

传统的人体姿态估计方法大多采用判别式（Discriminative）范式，即将姿态重构视为一个从传感器输入空间到关节坐标空间的静态映射问题。在这种框架下，模型仅关注如何最大化后验概率 $P(Y_t | X_t)$，即给定当前时刻的传感器观测 $X_t$，求解当前时刻的最优姿态 $Y_t$。然而，这种单向映射忽略了人体运动在时间维度上的连续性与因果依赖，导致模型本质上是一个静态的“估计器”而非动态的“预测器”。当面临需要预判动作趋势的场景时，此类方法无法利用历史运动规律进行有效推演，这正是现有研究缺乏动态预测能力的症结所在。为了突破这一局限，本研究引入自回归（Autoregressive, AR）生成机制，将人体姿态重构建模为基于时序历史的序列生成问题。

自回归生成模型的核心假设是：序列数据中某一时刻的状态严格依赖于其先前时刻的历史状态。对于人体运动序列 $S = \{s_1, s_2, ..., s_T\}$（其中 $s_t$ 可代表时刻 $t$ 的姿态向量或虚拟传感器数据），其联合概率分布 $P(S)$ 根据概率链式法则可分解为一系列条件概率的乘积：
\begin{equation}
P(S) = \prod_{t=1}^{T} P(s_t | s_1, s_2, ..., s_{t-1}) = \prod_{t=1}^{T} P(s_t | s_{<t})
\label{eq:ar_probability}
\end{equation}
该公式揭示了自回归机制的数学本质：当前时刻的输出不仅仅是当前输入的函数，更是整个历史轨迹的函数。在深度学习框架下，我们通常使用参数化神经网络（如前文所述的Causal Transformer）来拟合上述条件概率分布。这就要求模型在结构上必须遵循“因果约束”（Causal Constraint），即在计算时刻 $t$ 的特征时，严禁泄露时刻 $t$ 之后的未来信息。这通常通过在注意力机制中引入因果掩码（Causal Mask）来实现，确保感受野严格限制在历史时间窗内。

在推理阶段，自回归模型采用递归生成（Recursive Generation）的策略。不同于训练阶段可以使用真实标签进行并行计算（即Teacher Forcing模式），推理时模型需要将上一时刻生成的预测值 $\hat{s}_{t-1}$ 重新反馈至输入端，作为推演时刻 $t$ 的上下文依据。这种反馈回路（Feedback Loop）使得模型能够捕捉人体运动的内在动力学规律（Intrinsic Dynamics）。例如，在步态周期中，下肢的摆动具有高度的可预测性，模型可以通过学习历史时刻的角速度趋势，利用自回归机制在缺乏当前观测的情况下，“惯性”地推演出下一时刻的肢体位置。

具体到本研究的姿态重构任务中，自回归机制被用于构建IMU数据的预测闭环。模型不仅学习从IMU到姿态的映射，还同步学习IMU信号本身的时序演化规律：
\begin{equation}
\hat{x}_{t+1} = f_\theta(x_1, x_2, ..., x_t)
\label{eq:ar_prediction}
\end{equation}
其中 $f_\theta$ 为神经网络模型。通过不断将预测生成的虚拟IMU数据 $\hat{x}_{t+1}$ 注入输入流，系统能够利用历史信息构建未来的运动流形。这种生成式范式将姿态重构任务的定义从单纯的“状态恢复”拓展为“趋势预测”，从根本上解决了现有方法仅能感知当下而缺乏动态预测能力这一核心缺陷。

\xsection{本章小结}{Figures}

本章系统梳理了本文研究任务所涉及的核心理论与技术，为后续的方法设计与实验分析奠定了基础。内容首先涵盖了两类关键传感器：惯性测量单元与WiFi信道状态信息，明确了异构数据的物理感知特性；其次，探讨了Transformer与Mamba深度学习骨干网络，分析了其在长序列建模中的适用场景。在此基础上，重点阐述了多模态特征融合与对齐的关键技术，包括跨模态对比学习与交叉注意力机制，并介绍了用于实现动态时序推演的自回归生成理论。上述理论体系共同构成了解决异构数据深度融合及人体姿态动态预测问题的技术路径。